% ----------------------------------------------------------
% Capítulo de Resultados
% ----------------------------------------------------------
\chapter{Testes, validação e avaliação proposta}
\label{cap:resultados}
% ---

Este capítulo apresenta o estágio atual da validação do sistema. A comparação quantitativa com grandes modelos de linguagem ainda não foi executada, pois o sistema segue em aperfeiçoamento. Portanto, os resultados apresentados nesta etapa se concentram nos artefatos desenvolvidos, nos testes automatizados, nas verificações realizadas com o fluxo real de recuperação e no desenho da avaliação comparativa.

\section{Artefatos desenvolvidos}

O desenvolvimento resultou em uma aplicação funcional de consulta a manuais automotivos. Os principais artefatos são:

\begin{itemize}
  \item base local de documentos automotivos;
  \item pipeline de processamento de PDFs e geração de chunks;
  \item índice FAISS e metadados em JSON;
  \item script Python para recuperação híbrida;
  \item script Python para análise de imagens com CLIP;
  \item backend NestJS com autenticação, chats e orquestração RAG;
  \item frontend Next.js com streaming, fontes e galeria de imagens aprovadas;
  \item testes automatizados para reescrita de consultas, resposta e imagens.
\end{itemize}

\section{Testes automatizados}

Foram adicionados testes para cobrir falhas observadas durante o desenvolvimento. Os testes de reescrita de consultas verificam se perguntas incompletas são resolvidas corretamente quando o usuário envia uma mensagem curta de complemento. Também há testes para garantir que o retriever recebe a consulta resolvida, e não a mensagem bruta.

Os testes do serviço RAG cobrem casos como:

\begin{itemize}
  \item pergunta incompleta seguida de ano do veículo;
  \item uso de \texttt{resolvedQuery} no lugar de \texttt{rawMessage};
  \item substituição de resposta genérica por resposta operacional;
  \item rejeição de resposta que apenas orienta consultar o manual;
  \item correção de resposta que confunde tampa de óleo com reservatório;
  \item ausência de imagem confiável sem exibição de galeria.
\end{itemize}

Também foram criados testes específicos para o reranqueamento visual. Esses testes verificam a rejeição de capa de manual, imagem vazia, imagem genérica do veículo e figura de abertura do capô quando a pergunta é sobre a tampa de abastecimento do óleo.

\section{Validação com caso real}

Um caso usado durante a validação foi a pergunta: ``Onde fica a tampa do óleo do Audi A4 2008?''. Antes dos ajustes, o sistema podia recuperar uma figura relacionada à abertura do capô, que não respondia à pergunta. Após a alteração no pipeline, a recuperação textual passou a priorizar o trecho do manual do Audi A4 2008 referente à Figura 252, descrita como localização típica dos reservatórios, vareta de óleo e tampa de abastecimento do óleo do motor.

Na verificação manual do pipeline, o script de recuperação retornou como primeiro resultado o chunk do documento \texttt{2008-audi-a4-audi-a4\_2008\_bbb9f770c8cb74dba7ce.pdf}, páginas 257--258, contendo a expressão \textit{engine oil filler cap}. Também foi testado o CLIP com a imagem correta do compartimento do motor e uma imagem incorreta de abertura do capô. A mudança nos prompts positivos e negativos fez com que a imagem do compartimento do motor fosse priorizada.

Esse caso não representa ainda uma avaliação estatística, mas serviu para orientar melhorias importantes no fluxo: não dar bônus excessivo a qualquer imagem associada ao chunk, rejeitar figuras visualmente próximas porém incorretas e construir respostas textuais mais claras quando a evidência do manual é suficiente.

\section{Logs de depuração}

O backend registra logs estruturados com dados como mensagem original, pergunta pendente, consulta resolvida, páginas recuperadas, imagens selecionadas, imagens rejeitadas, razões de rejeição, confiança antes das penalizações e confiança final. Esses logs ajudam a investigar erros em conversas reais e também podem ser usados como base para futuras análises experimentais.

\section{Critérios para avaliação comparativa}

A avaliação comparativa proposta deverá utilizar um conjunto de perguntas representativas do uso real. As perguntas podem ser organizadas nas seguintes categorias:

\begin{itemize}
  \item especificações técnicas, como óleo, viscosidade, pneus e fluidos;
  \item localização de componentes, como tampas, reservatórios, fusíveis e bateria;
  \item interpretação de luzes e símbolos do painel;
  \item procedimentos de verificação e manutenção descritos no manual;
  \item perguntas ambíguas que exigem marca, modelo e ano;
  \item perguntas fora do escopo do manual.
\end{itemize}

Cada pergunta deverá possuir uma resposta de referência construída a partir do manual, incluindo documento, página e trecho utilizado. Dessa forma, será possível comparar as respostas do sistema RAG com as respostas de grandes modelos de linguagem sem acesso ao manual.

\section{Métricas propostas}

Para recuperação documental, pretende-se avaliar se o chunk correto aparece entre os primeiros resultados, a adequação da página recuperada, a compatibilidade do documento com o veículo e a utilidade do contexto retornado. Para a resposta final, os critérios propostos são correção factual, objetividade, completude, citação de fonte, ausência de invenção e utilidade prática.

Para perguntas visuais, a avaliação deverá verificar se a imagem exibida realmente ajuda a responder à pergunta. Uma imagem deve ser considerada útil quando mostra a peça, painel, reservatório, componente, seta, item numerado ou legenda correspondente. Imagens de capa, fotos genéricas do carro, figuras de outro procedimento ou imagens vazias devem ser consideradas incorretas.

\section{Desenho da comparação com grandes modelos}

A comparação planejada terá pelo menos três condições:

\begin{enumerate}
  \item sistema local com RAG e acesso aos manuais;
  \item modelo local sem recuperação documental;
  \item grandes modelos de linguagem sem acesso aos manuais.
\end{enumerate}

O objetivo será observar em quais tipos de pergunta o acesso ao manual produz ganho. A expectativa é que o RAG seja mais relevante em perguntas específicas e dependentes de documentação, enquanto modelos maiores podem manter vantagem em explicações gerais ou perguntas de conhecimento amplo.

\section{Limitações da validação atual}

A validação atual ainda é parcial. Os testes automatizados cobrem casos importantes, mas o conjunto de perguntas ainda precisa ser ampliado. A associação entre imagens e trechos de texto melhorou, porém ainda depende da qualidade dos PDFs e da extração das figuras. O cálculo de confiança também exige calibração com uma base maior de exemplos.

Mesmo assim, a etapa atual foi importante para estabilizar o fluxo conversacional e visual. O sistema passou a preservar melhor a intenção do usuário, a evitar imagens não aprovadas e a produzir respostas mais alinhadas aos trechos recuperados dos manuais.
